\documentclass[11pt]{article}

\usepackage[margin=1in]{geometry}
\usepackage[T1]{fontenc}
\usepackage[utf8]{inputenc}
\usepackage{lmodern}
\usepackage{microtype}
\usepackage{amsmath,amssymb,amsthm,mathtools}
\usepackage{bm}
\usepackage{enumitem}
\usepackage{hyperref}

\hypersetup{
    colorlinks=true,
    linkcolor=blue,
    citecolor=blue,
    urlcolor=blue
}

\newcommand{\R}{\mathbb{R}}
\newcommand{\E}{\mathbb{E}}
\newcommand{\Q}{\mathbb{Q}}
\newcommand{\cP}{\mathcal{P}}
\newcommand{\XiSet}{\Xi}
\newcommand{\eps}{\varepsilon}
\newcommand{\hxi}{\widehat\xi}
\newcommand{\hP}{\widehat{\mathbb P}_N}
\newcommand{\argmin}{\operatorname*{arg\,min}}
\newcommand{\argmax}{\operatorname*{arg\,max}}

\theoremstyle{plain}
\newtheorem{proposition}{Proposition}
\newtheorem{lemma}{Lemma}

\theoremstyle{definition}
\newtheorem{definition}{Definition}
\newtheorem{assumption}{Assumption}
\newtheorem{remark}{Remark}

\title{Wasserstein Distributionally Robust Nomination in Energy Markets}
\author{Charles Garrisi \qquad Welington de Oliveira \qquad François Pacaud}
\date{\today}

\begin{document}
\maketitle

\begin{abstract}
We consider the day-ahead nomination problem of a renewable producer exposed to a two-price imbalance settlement, in markets where prices may go negative.
Wasserstein distributionally robust optimization is a natural fit for this setting, where nominations must be committed with limited information.
However, the producer's loss couples uncertain generation with uncertain prices through bilinear terms, so that standard convex reformulations do not apply. 
We propose an exact solution method exploiting the low-dimensionality and the structure of the problem to reduce the search to a finite number of candidate points explicitly.
We compare our approach to conservative convex relaxations commonly employed.
Our experiments on synthetic and real-world data quantify the conservatism of the relaxations and show that exactness matters most when historical data is limited and when prices can go negative.
\end{abstract}

\vspace{2em}
\tableofcontents

\section{Introduction}
\subsection{Motivation}
\subsection{Literature review}
\subsubsection{Classical methods are not enough (SO/SAA, RO)}
\subsubsection{DRO is the way to go}
\subsection{Contributions}
\subsection{Paper structure}


\section{Problem formulation}
We consider a renewable-energy producer participating in a day-ahead electricity market and subsequently exposed to an imbalance-settlement mechanism. 
The decision is made before delivery. 
At that time, the producer chooses a nomination $n$, the decision variable representing the quantity of energy committed to the day-ahead market, while the realized generation $g$ and the relevant market prices are still uncertain.
These uncertain parameters are denoted by the vector $\xi$.
Generation and nominations are normalized by the installed capacity, so that the feasible decision set for $n$ is $\mathcal{X}=[0,1]$.
Lastly, we focus on a single delivery period — which corresponds to a 15-minute interval for the European market since October $1^{\text{st}}$, 2025 — and adopt the convention that the producer aims to minimize a loss function denoted by $\ell(n,\xi)$. 

\subsection{Uncertainty modeling} 
\label{subsec:uncertainty-modeling} 

\subsubsection{Description of the uncertainty structure}
As mentioned, we model the random environment by a vector $\xi$ of uncertain parameters. 
This vector contains the realized renewable generation $g$, the day-ahead spot price $s$, the up-regulation price $r^+$, and the down-regulation price $r^-$. 
The variable $g\in[0,1]$ denotes the actual production as a realized fraction of the available capacity. 
The spot price $s$ is the price at which the day-ahead nomination is settled, while the prices $r^+$ and $r^-$ determine the cost or remuneration of deviations between the nomination and the realized generation. 
If the producer is short, meaning $n>g$, the deficit $(n-g)^+$ is settled at the up-regulation price $r^+$. 
If the producer is long, meaning $g>n$, the surplus $(g-n)^+$ is settled at the down-regulation price $r^-$. 
The single-period loss (i.e. the negative profit) is therefore expressed as 
\begin{equation} 
    \label{eq:loss-original} 
    \ell(n,\xi) = -ns - r^-(g-n)^+ + r^+(n-g)^+ . 
\end{equation} 
The first term corresponds to the day-ahead revenue while the last two terms encode the imbalance settlement. 
We do not require the spot price $s$ or the imbalance prices $r^+, r^-$ to be nonnegative.
Negative prices are now a structural feature of electricity markets.
The day-ahead price was negative for 513 hours in 2025 (about $6\%$ of the year) up from 352 in 2024 and 147 in 2023 \color{red}[\href{https://analysesetdonnees.rte-france.com/bilan-electrique-2025/prix#Heuresaprixnegatif}{source}]\color{black}. 
These negative hours tend to cluster at midday between April and August, when solar and wind output is high and demand is low.
French TSO RTE expects their frequency to keep rising with renewable penetration.
In the first half of 2026, entire days with a negative average price appear in the data.
This is not an isolated phenomenon: negative prices have also been observed in Germany, Belgium, Spain, and Switzerland.
\\

\noindent Allowing negative prices completely alters the producer's economic incentives and hence changes the decision. 
For instance, a revenue term in~\eqref{eq:loss-original} is decreasing in $n$ when $s>0$ (larger nominations are rewarded) but increasing when $s<0$ (larger nominations are penalized).
A negative spot price can thus flip the marginal value of committing energy. 
Negative prices also affect the geometry of the uncertainty set and invalidate convexity properties that hold in classical positive-price newsvendor models. 
\\

\subsubsection{How to model generation?}
\label{subsubsec:model-generation}

The realized renewable production $g$ is not known at the time of nomination. 
In practice, the producer observes a point forecast, or possibly a predictive distribution, but the realized production may differ because of weather forecast errors, local curtailment, outages, or measurement noise. 
Since both the nomination and the realized production are normalized by the available capacity, 
\[
    g\in[\underline g,\overline g]\subseteq[0,1].
\]
The bounds $\underline g$ and $\overline g$ may be interpreted in different ways depending on the application. 
They may correspond to physical limits, to forecast-based confidence bounds, or to empirical bounds obtained after removing gross data errors. 
In the numerical experiments, we use the normalized bounds $\underline g=0$ and $\overline g=1$.
\\

\noindent It is useful to think of realized generation as a perturbation of a forecast. 
Let $f\in[0,1]$ denote a point forecast and let $\zeta$ denote a relative forecast error, so that the realized generation becomes
\[
    g=\Pi_{[0,1]}\big((1+\zeta)f\big),
\]
where $\Pi_{[0,1]}$ denotes projection onto the interval $[0,1]$. 
This representation has two advantages. 
First, it preserves the interpretation of $g$ as a normalized production level. 
Second, it makes forecast errors scale with the forecast itself: an error of a given relative magnitude has a larger absolute impact when expected production is high.
\\

\noindent This parametric representation is only used to generate synthetic data in the numerical experiments. 
The optimization model itself does not assume anything on $g$ apart from its support bounds. 


\subsubsection{How to model spot prices?}
\label{subsubsec:model-spot-price}

The day-ahead spot price $s$ is also uncertain at the time of nomination. 
It is determined by the market-clearing process after participants submit their bids \color{red}[At EDF summer school, someone suggested to add these aggregated bids and especially the "best bid" as an additional uncertainty variable, could be interesting]\color{black}, reflecting the interaction between demand, renewable availability, generation costs, plant availability, interconnection constraints, and network conditions. 
Spot prices are therefore hard to predict by essence and volatile.
They may exhibit spikes, strong seasonality, and regime changes.
\\

\noindent As mentioned, a central feature of modern electricity markets is that spot prices may become negative. 
Negative prices typically occur when inflexible, high renewable availability, and low demand jointly create excess supply.
For this reason, we do not impose positive spot prices. 
Instead, we assume that $s$ is confined within a box;, so that $s\in[\underline s,\overline s]$, where $\underline s$ may be negative.
\\

\noindent In the synthetic experiments, spot prices are generated from a rescaled Student's t-distribution \color{red}[This is to discuss] \color{black}so that they can take negative values and exhibit heavier tails than a normal distribution.
This choice is not part of the model itself but provides a controlled way to test the sensitivity of the different methods to volatility, price spikes, support size, and negative-price regimes.

\subsubsection{How to model imbalance prices?}
\label{subsubsec:model-imbalance-price}

Imbalance prices are knowingly very difficult to model.
Rather than modelling $r^+$ and $r^-$ directly, we separate the average imbalance level from the imbalance asymmetry. 
Let \[ r=\frac{r^++r^-}{2}, \qquad \delta=r-s, \] where $\delta$ measures the deviation between the average imbalance price and the spot price. 
The variable $\delta$ captures the residual imbalance component that is not explained by the spot price alone. 
Positive values of $\delta$ correspond to an average imbalance price above the spot price, while negative values correspond to an average imbalance price below the spot price.
We then introduce an asymmetry parameter $\tau\in[0,1)$ and write
\begin{equation}
    \label{eq:imbalance-reparam}
    r^+=(s+\delta)(1+\tau),
    \qquad
    r^-=(s+\delta)(1-\tau),
\end{equation}
where $\tau$ controls the relative spread between the two imbalance prices around their average level $r=s+\delta$. 
This parametrization separates the stochastic component $\delta$, which measures the average imbalance level relative to the spot price, from the structural asymmetry $\tau$, which controls the difference between deficit and surplus settlement.
\\

\noindent In the following, $\tau$ is assumed to be fixed over a given delivery period.
This enables to keep the asymmetry of the two-price imbalance mechanism explicit while avoiding an additional random component. 
From the data, this assumption appears reasonable. \color{red}[Figure to prove my point] \color{black}
After this reparametrization, the uncertainties vector becomes
\[
    \xi=(g,s,\delta)\in\R^3.
\]
The corresponding support is then
\[
    \XiSet
    =
    [\underline g,\overline g]
    \times
    [\underline s,\overline s]
    \times
    [\underline\delta,\overline\delta].
\]
The imbalance spreads faced by the producer are
\begin{equation}\label{eq:marginal_cost}
    c^-(s,\delta)
    =
    s-r^-
    =
    (s+\delta)\tau-\delta,
    \qquad
    c^+(s,\delta)
    =
    r^+-s
    =
    (s+\delta)\tau+\delta .
\end{equation}
Here, in~\eqref{eq:marginal_cost}, $c^-(s,\delta)$ is the marginal loss associated with a production surplus, while $c^+(s,\delta)$ is the marginal loss associated with a production deficit. 
Both depend on uncertain prices. 
This is a major difference from the classical newsvendor model, where the underage and overage costs are deterministic constants.
Using the identity
\[
    n-g=(n-g)^+-(g-n)^+,
\]
the loss can be rewritten as
\begin{equation}
    \label{eq:loss-reparam}
    \ell(n,\xi)
    =
    -gs
    +
    c^-(s,\delta)(g-n)^+
    +
    c^+(s,\delta)(n-g)^+,
    \qquad
    \xi=(g,s,\delta).
\end{equation}
The term first term$-gs$ in~\eqref{eq:loss-reparam} is the negative value of realized production at the spot price, while the remaining terms penalize deviations between realized production and nomination, with imbalance costs that depend on the market state. 
This formulation makes explicit the newsvendor-like structure of the nomination decision, but it also reveals the main difficulty: the objective contains products of uncertainty variables, in particular $gs$ and $g\delta$. 
Consequently, while the loss is piecewise affine in the nomination for fixed uncertainty, it is nonconvex in the uncertainties. 
\\

\noindent In the synthetic experiments, imbalance spreads are generated from a combination of two light-tailed distributions \color{red}[This is to discuss] \color{black}, corresponding to a normal and an extreme market regimes, and then truncated to the support interval. 
This choice is not part of the model itself but provides a controlled way to test the sensitivity of the different methods to volatility, price spikes, support size, and negative-price regimes.
\color{red}
\begin{itemize}
    \item Check if having $r=s+\delta<0$ is a problem or not
    \item Check if having $c^+ + c^-$ negative is a problem or not (non-convexity?)
\end{itemize}
\color{black}
\subsection{Optimization problem} 
\label{subsec:optimization-problem} 

\subsubsection{Assumptions}

We now state the modelling assumptions used throughout the paper. 
\color{red}
We should guarantee that
\begin{equation}
    \mathbb{E}[r^-] < \mathbb{E}[s] < \mathbb{E}[r^+],
\end{equation}
otherwise the producer would always be better off being long or short. 
\color{black}
\begin{assumption}[Price-taking producer] 
    \label{ass:price-taker} The producer is a price taker. 
    We consider a producer small enough so that its nomination does not affect the spot price, the imbalance prices, or the system imbalance state. 
\end{assumption}
Without this assumption, the nomination would influence the spot or imbalance prices.
Except for big actors like EDF in France, this is generally the case.

\begin{assumption}[Box support] 
    \label{ass:compact-support} 
    The uncertainties vector satisfies \[ \xi=(g,s,\delta)\in\XiSet := [\underline g,\overline g] \times [\underline s,\overline s] \times [\underline\delta,\overline\delta], \] where $0\le \underline g\le \overline g\le1$. 
    This is more constraining than the usual compactness assumption.
    Note that the price bounds may contain negative values. 
\end{assumption} 
This assumption ensures that the loss is bounded and attains its worst-case value on $\XiSet$. 
This is needed for the sample-wise separation problems to be well posed and for the finite-candidate arguments used later: over a box support, the bilinear terms can be lifted, bounded, and relaxed by McCormick envelopes, and the exact separation problem can be reduced to a finite set of candidate points. 

\begin{assumption}[Historical data] 
    \label{ass:data} 
    We observe $N$ historical samples \[ \hxi_i=(\widehat g_i,\widehat s_i,\widehat\delta_i)\in\XiSet, \qquad i=1,\ldots,N, \] and denote the empirical distribution by \[ \hP=\frac1N\sum_{i=1}^N\delta_{\hxi_i}. \] 
\end{assumption} 
The empirical distribution is the center of the Wasserstein ambiguity set as this is our best known estimate at the time of the decision. 

\begin{assumption}[No independence] 
    \label{ass:no-independence} 
    We impose no independence between $g$, $s$, and $\delta$. 
    The Wasserstein ambiguity set is built on the joint empirical distribution of $\xi=(g,s,\delta)$. 
\end{assumption} 

Renewable generation, spot prices, and imbalance prices are typically dependent: high renewable availability often coincide with low or negative spot prices, and imbalance prices react to system conditions that are themselves correlated with renewable forecast errors. 


\begin{assumption}[Transportation metric] 
    \label{ass:transportation-metric} 
    The transportation cost is the weighted $\ell_1$ metric \[ d_W(\xi,\xi') = \|W(\xi-\xi')\|_1 = w_{\xi_1}|\xi_1-\xi_1'|+\dots+w_{\xi_d}|\xi_d-\xi_d'|, \] where $\xi\in\mathbb{R}^d$, $W=\operatorname{diag}(w_{\xi_1},\dots,w_{\xi_d})$, and $w_{\xi_1},\dots,w_{\xi_d}>0$. 
\end{assumption} 

The weights in $W$ normalize the different physical units and magnitudes of the uncertainty components. 
Without such a normalization, the Wasserstein ball would be dominated by the component with the largest numerical scale, which would completely distort the ambiguity geometry. 
In the numerical experiments, these weights are computed online based on the empirical standard deviations of the recent observations. 
Another approach is to use the median absolute deviation (MAD) or the interquartile range (IQR) that have the advantage of being more robust estimators.

\subsubsection{DRO formulation}
Let us write the distributionally robust formulation for our nomination problem.
Given a Wasserstein radius $\eps\ge0$, we define the ambiguity set 
    \begin{equation} 
        \label{eq:wasserstein-ambiguity} 
        \cP_\eps = \left\{ \Q\in\cP(\XiSet): \mathcal{W}_{d_W}(\Q,\hP)\le \eps \right\}, 
    \end{equation} 
    where $\mathcal{W}_{d_W}$ denotes the Wasserstein distance induced by the ground metric $d_W$. 
    The resulting distributionally robust nomination problem is 
    \begin{equation} \label{eq:dro-primal-main} \inf_{n\in[0,1]} \sup_{\Q\in\cP_\eps} \E_\Q[\ell(n,\xi)] . \end{equation} 
    The radius $\eps$ controls the level of robustness. When $\eps=0$, the ambiguity set reduces to the empirical distribution and \eqref{eq:dro-primal-main} becomes the sample average approximation \begin{equation} \label{eq:saa-pb} \min_{n\in[0,1]} \frac1N\sum_{i=1}^N \ell(n,\hxi_i). \end{equation}
    As $\eps$ increases, the adversary can move probability mass further away from the observed samples, producing decisions that are more conservative but less sensitive to sampling error. 
    In the limiting regime where the transportation budget is sufficiently large, the behavior approaches that of a robust model driven primarily by the support $\XiSet$. 
    For fixed $\xi$, the function $n\mapsto\ell(n,\xi)$ is piecewise affine. 
    It is convex in $n$ whenever \begin{equation} \label{eq:convexity-condition} c^-(s,\delta)+c^+(s,\delta) = 2\tau(s+\delta) \ge0. \end{equation}
    This condition is met when $\tau\ge0$ and $s+\delta=r\ge0$, but it may fail if the average imbalance price becomes negative. 
    For fixed $n$, the map $\xi\mapsto\ell(n,\xi)$ is nonconvex because of the products $gs$ and $g\delta$. 
    This prevents a direct use of classical convex reformulations in the likes of~\cite{EsfahaniKuhn2018} that rely on convexity of the loss in the uncertain parameters. 
    The difficulty is now to solve~\eqref{eq:dro-primal-main} which is a nonconvex infinite-dimensional moment problem.
    
\section{Solution methods} 
\label{sec:solution-methods} 
This section presents three resolution strategies.
The first two are conservative convex reformulations based on lifting the uncertainty variables and relaxing the support set to bypass the nonconvexity. 
The third one keeps the nonconvexity and exploits both the structure and the low-dimensionality of the problem. 

 \subsection{Relaxation as a convex problem} 
 \label{subsec:convex-relaxations} 
 \subsubsection{Lifting the bilinear terms}
 The nonconvexity in $\xi$ arises from the products $gs$ and $g\delta$ in the loss function~\eqref{eq:loss-reparam}. 
 To address this, we introduce lifted variables \[ t_s=gs, \qquad t_\delta=g\delta, \] and define the lifted uncertainties vector \[ \widetilde\xi=(g,s,\delta,t_s,t_\delta)\in\R^5. \] 
 In the lifted space, the two branches of the loss become affine in $\widetilde\xi$. 
 On the surplus branch $g\ge n$, one has \begin{align} \label{eq:loss-surplus-lifted} \ell^+(n,\widetilde\xi) &= -t_s+\big((s+\delta)\tau-\delta\big)(g-n) \nonumber\\ &= (\tau-1)t_s + (\tau-1)t_\delta - \tau n s + (1-\tau)n\delta . \end{align} 
 On the deficit branch $g\le n$, one has \begin{align} \label{eq:loss-deficit-lifted} \ell^-(n,\widetilde\xi) &= -t_s+\big((s+\delta)\tau+\delta\big)(n-g) \nonumber\\ &= -(1+\tau)t_s - (1+\tau)t_\delta + \tau n s + (1+\tau)n\delta . \end{align} 
 For fixed $n$, both branches can be written as \[ \ell^k(n,\widetilde\xi)=a_k(n)^\top\widetilde\xi+b_k(n), \qquad k\in\{+,-\}, \] with $b_+(n)=b_-(n)=0$ and \[ a_+(n) = \begin{pmatrix} 0\\ -\tau n\\ (1-\tau)n\\ \tau-1\\ \tau-1 \end{pmatrix}, \qquad a_-(n) = \begin{pmatrix} 0\\ \tau n\\ (1+\tau)n\\ -(1+\tau)\\ -(1+\tau) \end{pmatrix}. \] 
 The exact lifted support is \[ \widetilde\Xi_{\rm exact} = \left\{ (g,s,\delta,t_s,t_\delta): (g,s,\delta)\in\XiSet,\ t_s=gs,\ t_\delta=g\delta \right\}. \] 
 However, this set is still nonconvex. 
 The following relaxations replace $\widetilde\Xi_{\rm exact}$ by a tractable polyhedral outer approximation. 
 
 \subsubsection{Linear relaxations}
\paragraph{Box linearization.}
 The simplest relaxation consists in keeping only coordinate-wise bounds on the lifted variables: \[ \widetilde\Xi_{\rm box} = [\underline g,\overline g] \times [\underline s,\overline s] \times [\underline\delta,\overline\delta] \times [\underline t_s,\overline t_s] \times [\underline t_\delta,\overline t_\delta], \] where \[ \underline t_s = \min\{gs:g\in[\underline g,\overline g],\ s\in[\underline s,\overline s]\}, \qquad \overline t_s = \max\{gs:g\in[\underline g,\overline g],\ s\in[\underline s,\overline s]\}, \] and similarly \[ \underline t_\delta = \min\{g\delta:g\in[\underline g,\overline g],\ \delta\in[\underline\delta,\overline\delta]\}, \qquad \overline t_\delta = \max\{g\delta:g\in[\underline g,\overline g],\ \delta\in[\underline\delta,\overline\delta]\}. \] 
 Equivalently, \[ \widetilde\Xi_{\rm box} = \{\widetilde\xi\in\R^5:C_{\rm box}\widetilde\xi\le d_{\rm box}\}. \] 
 This relaxation is very cheap computationally, but it is generally very loose because it ignores the dependence between $g$, $s$, $\delta$, and the lifted products. 
 \\

More than that, the fact that the lifted products are decoupled from the original variables entails that this approach may lead to impossible scenarios in the lifted space. \color{red}[Important remark, insist on this?] \color{black}
\color{red}[Here, maybe add the formal proof or an example of this] \color{black}
In particular, as $s$ and $\delta$ can be negative while $g>0$, the box relaxation includes scenarios with $t_s=gs<0$ or $t_\delta=g\delta<0$ even though $g$ and $s$ (respectively $\delta$) are both nonnegative in the original space. 
This problem is discussed in more detail in the numerical experiments, where we observe that the box relaxation is so loose that it can lead to extremely conservative decisions — essentially a zero nomination.

\paragraph{McCormick linearization.}
We can obtain a tighter convex relaxation by replacing the bilinear equalities by their McCormick envelopes. 
As every uncertainty variable is confined within a box, we can write explicit linear inequalities and partially restore the dependence between $g$, $s$, and $\delta$.
\\

\noindent For instance, for $t_s=gs$, we know that $g\in[\underline g,\overline g]$ and $s\in[\underline s,\overline s]$, so that $(g-\underline g)(s-\underline s)\ge0$ and its envelope is given by 
\begin{align} t_s &\ge \underline g\,s+\underline s\,g-\underline g\,\underline s, \label{eq:mccormick-s-1}\\ t_s &\ge \overline g\,s+\overline s\,g-\overline g\,\overline s, \label{eq:mccormick-s-2}\\ t_s &\le \underline g\,s+\overline s\,g-\underline g\,\overline s, \label{eq:mccormick-s-3}\\ t_s &\le \overline g\,s+\underline s\,g-\overline g\,\underline s. \label{eq:mccormick-s-4} \end{align} 
Analogously, for $t_\delta=g\delta$, \begin{align} t_\delta &\ge \underline g\,\delta+\underline\delta\,g-\underline g\,\underline\delta, \label{eq:mccormick-d-1}\\ t_\delta &\ge \overline g\,\delta+\overline\delta\,g-\overline g\,\overline\delta, \label{eq:mccormick-d-2}\\ t_\delta &\le \underline g\,\delta+\overline\delta\,g-\underline g\,\overline\delta, \label{eq:mccormick-d-3}\\ t_\delta &\le \overline g\,\delta+\underline\delta\,g-\overline g\,\underline\delta. \label{eq:mccormick-d-4} \end{align} 
Compared to the box linearization, this approach provides a tighter relaxation, although it is of course still looser than the exact formulation. 
Let \[ \widetilde\Xi_{\rm Mc} = \{\widetilde\xi\in\R^5:C_{\rm Mc}\widetilde\xi\le d_{\rm Mc}\} \] denote the polytope defined by the original bounds and by \eqref{eq:mccormick-s-1}--\eqref{eq:mccormick-d-4}. 
So far, we have that \[ \widetilde\Xi_{\rm exact} \subseteq \widetilde\Xi_{\rm Mc} \subseteq \widetilde\Xi_{\rm box}. \] 
Consequently, for every fixed $n$, $\lambda$, and sample $i$, the separation value is bounded as follows: 
\[ \mathcal V_i^{\rm exact}(n,\lambda) \le \mathcal V_i^{\rm Mc}(n,\lambda) \le \mathcal V_i^{\rm box}(n,\lambda). \] 

\subsubsection{Wasserstein DRO and the Esfahani--Kuhn LP reformulation}
\label{subsubsec:ek-lp-reformulation}

We now show how these lifted relaxations lead to tractable Wasserstein DRO problems. 
The key point is that, after lifting the bilinear terms and replacing the lifted support by a polyhedral outer approximation, the relaxed loss can be written as the maximum of finitely many affine functions of the lifted uncertainties vector. 
Indeed, we are now in the setting in which the convex reformulation theorem of Esfahani and Kuhn~\cite{EsfahaniKuhn2018} yields a finite-dimensional LP.
We can rewrite the lifted loss as
\begin{equation}\label{eq:max-affine-loss}
    \widetilde\ell_R(n,\widetilde\xi)
    =
    \max_{k\in\{+,-\}}
    \left\{
        a_k(n)^\top\widetilde\xi+b_k(n)
    \right\}.
\end{equation}
To simplify the notation in the following, we will simply use $\widetilde\ell_R(n,\xi)$ to denote the lifted loss.

\begin{proposition}[Esfahani--Kuhn convex reformulation]
\label{prop:convex-reformulation-kuhn}
Suppose that the uncertainty set $\Xi\subseteq\mathbb{R}^m$ is convex and closed, and the negative constituent function $-\ell_k$ are proper, convex, and lower semicontinuous for all $k\le K$. 
Moreover, suppose that $\ell_k$ is not identically equal to $-\infty$ for all $k\le K$.
Then, for any $\varepsilon\ge 0$, the lifted Wasserstein DRO problem
\[
    \inf_{n\in[0,1]}
    \sup_{\mathbb Q\in{\mathcal P}_{\varepsilon}}
    \mathbb E_{\mathbb Q}
    \left[
        \widetilde\ell(n,\widetilde\xi)
    \right]
\]
is equivalent to the LP
\begin{equation}
\label{eq:lifted-dro-lp}
\begin{aligned}
    \min_{n,\lambda,z,\gamma}\quad
    &
    \lambda\varepsilon+\frac1N\sum_{i=1}^N z_i
    \\
    \textnormal{s.t.}\quad
    &
    b_k(n)
    +
    a_k(n)^\top\widehat{\xi}_i
    +
    \gamma_{ik}^\top
    \left(
        d-C\widehat{\xi}_i
    \right)
    \le z_i,
    &&
    i=1,\ldots,N,\ k\in\{+,-\},
    \\
    &
    C^\top\gamma_{ik}-a_k(n)\le \lambda\widetilde w,
    &&
    i=1,\ldots,N,\ k\in\{+,-\},
    \\
    &
    -C^\top\gamma_{ik}+a_k(n)\le \lambda\widetilde w,
    &&
    i=1,\ldots,N,\ k\in\{+,-\},
    \\
    &
    \gamma_{ik}\ge0,
    &&
    i=1,\ldots,N,\ k\in\{+,-\},
    \\
    &
    0\le n\le1,
    \qquad
    \lambda\ge0 .
\end{aligned}
\end{equation}
\end{proposition}

\begin{proof}Esfahani and Kuhn~\cite{EsfahaniKuhn2018}.\\

 Let $ W=\operatorname{diag}( w)$ be the weight matrix used in the lifted transportation cost, and let \[ \widehat{\xi}_i = (\widehat g_i,\widehat s_i,\widehat\delta_i, \widehat g_i\widehat s_i, \widehat g_i\widehat\delta_i) \] be the natural lifting of sample $i$. 
 Consider a generic polyhedral relaxation \[ \widetilde\Xi = \{\xi:C\xi\le d\}.\]
 For each branch $k\in\{+,-\}$, the relaxed separation term is \[ \sup_{\xi\in\widetilde\Xi} \left\{ a_k(n)^\top\xi+b_k(n) - \lambda \|\widetilde W(\xi-\widehat{\xi}_i)\|_1 \right\}. \] 
 By LP duality, this supremum is bounded above by $z_i$ if and only if there exists $\gamma_{ik}\ge0$ such that \[ C^\top\gamma_{ik}-a_k(n)\le \lambda \widetilde w, \qquad -C^\top\gamma_{ik}+a_k(n)\le \lambda \widetilde w, \] and \[ b_k(n) + a_k(n)^\top\widehat{\xi}_i + \gamma_{ik}^\top \left(d-C\widehat{\xi}_i\right) \le z_i . \] 
\end{proof}


\subsection{Exact reformulation} 
\label{subsec:exact-reformulation} 
We now describe an exact approach for evaluating the separation value
\begin{equation}\label{eq:separation-problem}
\mathcal V_i(n,\lambda) = \sup_{\xi\in\XiSet} \left\{ \ell(n,\xi)-\lambda d_W(\xi,\hxi_i) \right\}
\end{equation}
without relaxing the bilinear terms.
The key observation is that, given the low-dimensional structure of our problem, we can reduce the search from $\Xi$ to a finite set of candidate points, given by the KKT necessary conditions. 
That way, we can evaluate the separation problem both efficiently and exactly.

\begin{remark}
The convex relaxations of Section~\ref{subsec:convex-relaxations} operate on the lifted vector $(g,s,\delta,t_s,t_\delta)$, whereas the exact method stays in the original space. 
Yet, the two are still comparable because the transportation is measured on $(g,s,\delta)$ for every method.
To do so, the metric weights of the auxiliary variables $(t_s,t_\delta)$ are set to zero in the lifted formulations.
\end{remark}

\subsubsection{Exact finite reduction of the separation problem}
\label{subsubsec:finite-separation}
 
In this section, we fix a nomination $n\in[0,1]$, a dual multiplier $\lambda\ge0$, and a sample index $i\in\{1,\dots,N\}$. 
Since $\Xi=(\Xi\cap\{g\le n\})\cup(\Xi\cap\{g\ge n\})$, we may split the separation problem into the surplus and deficit branches: \[ \mathcal V_i(n,\lambda) = \max\{\mathcal V_i^+(n,\lambda),\mathcal V_i^-(n,\lambda)\}, \] where 
\begin{align} \label{eq:sep-plus} \mathcal V_i^+(n,\lambda) &:= \sup_{\xi\in\XiSet,\ g\ge n} \left\{ -gs+c^-(s,\delta)(g-n) -\lambda d_W(\xi,\hxi_i) \right\}, \\ \label{eq:sep-minus} \mathcal V_i^-(n,\lambda) &:= \sup_{\xi\in\XiSet,\ g\le n} \left\{ -gs+c^+(s,\delta)(n-g) -\lambda d_W(\xi,\hxi_i) \right\}, \end{align} 
with the convention that a supremum over an empty set equals $-\infty$.
To lighten notation we drop the index $i$ and write $\hxi=(\widehat g,\widehat s,\widehat\delta)\in\XiSet$ for the sample so that the objective in~\eqref{eq:separation-problem} becomes
\begin{equation}
    \label{eq:phi-def}
    \varphi(\xi)
    :=
    \ell(n,\xi)-\lambda\, d_W(\xi,\hxi),
    \qquad
    \xi=(g,s,\delta)\in\XiSet.
\end{equation}

\noindent The function $\varphi$ is continuous but neither convex nor concave as the loss \eqref{eq:loss-reparam} contains the bilinear terms $gs$ and $g\delta$.
The absolute-value terms introduce breakpoints at $g=\widehat g_i$, $s=\widehat s_i$, and $\delta=\widehat\delta_i$.
The loss also introduces a breakpoint at the nomination $g=n$ which is quite peculiar.
In practice, the generation almost surely differs from the nomination as the grid is never perfectly balanced, and the generation forecast from the actual generation, but $n$ is still included in the candidate set $G_i$.
\color{red}[Maybe add a numerical example to show that we can have $g=n$ as a valid theoretical worst-case?] \color{black}
We now show in Lemma~\ref{lemma:candidate-sets} that for each sample $i$ the supremum is attained on a finite set of candidate points.

\begin{lemma}[Exact finite reduction]
\label{lemma:candidate-sets}
Let Assumptions~\ref{ass:compact-support} and~\ref{ass:transportation-metric} hold, and fix $n\in[0,1]$, $\lambda\ge0$, and $i\in\{1,\dots,N\}$. Define the candidate sets
\[
    G_i(n)
    :=
    \{\underline g,\overline g,\widehat g_i\}
    \cup
    \bigl(\{n\}\cap[\underline g,\overline g]\bigr),
    \qquad
    S_i:=\{\underline s,\overline s,\widehat s_i\},
    \qquad
    D_i:=\{\underline\delta,\overline\delta,\widehat\delta_i\},
\]
and the candidate grid $\mathcal C_i(n):=G_i(n)\times S_i\times D_i\subseteq\XiSet$.
Then the supremum is attained and
\begin{align}
\label{eq:finite-separation}
\mathcal V_i(n,\lambda)
&= \max_{\xi\in\mathcal C_i(n)} \varphi(\xi)\notag\\
&=
\max_{\xi\in\mathcal C_i(n)}
\left\{
    \ell(n,\xi)-\lambda\,d_W(\xi,\hxi_i)
\right\} \notag\\
&=
\max\Bigg\{
    \max_{\substack{
        g\in G_i(n),\,s\in S_i,\,\delta\in D_i\\
        g\ge n
    }}
    \left[
        -gs+c^-(s,\delta)(g-n)
        -\lambda d_W\bigl((g,s,\delta),\hxi_i\bigr)
    \right],
    \notag\\
&\hspace{1.6cm}
    \max_{\substack{
        g\in G_i(n),\,s\in S_i,\,\delta\in D_i\\
        g\le n
    }}
    \left[
        -gs+c^+(s,\delta)(n-g)
        -\lambda d_W\bigl((g,s,\delta),\hxi_i\bigr)
    \right]
\Bigg\}.
\end{align}
\end{lemma}
 

\begin{proof}
We first establish that the supremum is attained on $\XiSet$.
We then study the deficit branch $g\le n$ fully by writing the complete KKT system and show in each case that a global maximizer can be chosen in the finite candidate set $\mathcal C_i(n)$.
The surplus branch is treated by the same argument.
Taking the maximum over the two branches concludes.
\\

\noindent The maps $\xi\mapsto(g-n)^+$, $\xi\mapsto(n-g)^+$, and $\xi\mapsto d_W(\xi,\hxi)$ are continuous, and $c^\pm(s,\delta)$ are affine, so that $\varphi$ is continuous on the nonempty compact set $\XiSet$. 
By the Weierstrass theorem the supremum in~\eqref{eq:separation-problem} is attained at some $\xi^\star=(g^\star,s^\star,\delta^\star)\in\XiSet$, with value $\varphi(\xi^\star)$. 
Since $\mathcal C_i(n)\subseteq\XiSet$, the right-hand side of \eqref{eq:finite-separation} is at most $\varphi(\xi^\star)$.
Both branches being analogous, we detail the deficit branch $(g\le n)$.
Introducing epigraph variables $u=(u_g,u_s,u_\delta)$ for the three absolute values, the separation problem for the deficit branch is equivalent to
\begin{equation}
\label{eq:deficit-epigraph}
\begin{aligned}
    \max_{g,s,\delta,u}
    \quad
    &
    -gs
    +
    \bigl[(s+\delta)\tau+\delta\bigr](n-g)
    -
    \lambda
    \left(
        w_g u_g
        +
        w_s u_s
        +
        w_\delta u_\delta
    \right)
    \\
    \textnormal{s.t.}\quad
    &
    g-n\le0,
    &&
    (\mu)
    \\
    &
    \underline g-g\le0,
    \qquad
    g-\overline g\le0,
    &&
    (\nu_1,\ \nu_2)
    \\
    &
    \underline s-s\le0,
    \qquad
    s-\overline s\le0,
    &&
    (\eta_1,\ \eta_2)
    \\
    &
    \underline\delta-\delta\le0,
    \qquad
    \delta-\overline\delta\le0,
    &&
    (\rho_1,\ \rho_2)
    \\
    &
    (g-\widehat g)-u_g\le0,
    \qquad
    (\widehat g-g)-u_g\le0,
    &&
    (\alpha_1,\ \alpha_2)
    \\
    &
    (s-\widehat s)-u_s\le0,
    \qquad
    (\widehat s-s)-u_s\le0,
    &&
    (\beta_1,\ \beta_2)
    \\
    &
    (\delta-\widehat\delta)-u_\delta\le0,
    \qquad
    (\widehat\delta-\delta)-u_\delta\le0.
    &&
    (\gamma_1,\ \gamma_2)
\end{aligned}
\end{equation}       

\noindent Since all constraints of \eqref{eq:deficit-epigraph} are affine, the Abadie constraint qualification hold and the Karush-Kuhn-Tucker conditions are necessary at $\xi^\star$.
Forming the Lagrangian, 
\begin{align*}
    \mathcal L
    ={}&
    -gs+\bigl[(s+\delta)\tau+\delta\bigr](n-g)
    -\lambda\left(w_g u_g+w_s u_s+w_\delta u_\delta\right)
    -\mu(g-n)
    \\
    &
    -\alpha_1\bigl(g-\widehat g-u_g\bigr)
    -\alpha_2\bigl(\widehat g-g-u_g\bigr)
    -\beta_1\bigl(s-\widehat s-u_s\bigr)
    -\beta_2\bigl(\widehat s-s-u_s\bigr)
    \\
    &
    -\gamma_1\bigl(\delta-\widehat\delta-u_\delta\bigr)
    -\gamma_2\bigl(\widehat\delta-\delta-u_\delta\bigr)
    -\nu_1(\underline g-g)
    -\nu_2(g-\overline g)
    \\
    &
    -\eta_1(\underline s-s)
    -\eta_2(s-\overline s)
    -\rho_1(\underline\delta-\delta)
    -\rho_2(\delta-\overline\delta),
\end{align*}
the stationarity conditions are
\begin{align}
    \frac{\partial\mathcal L}{\partial g}=0
    &\iff
    -(1+\tau)(s+\delta)
    -\mu-\alpha_1+\alpha_2+\nu_1-\nu_2=0,
    \label{eq:stat-g}
    \\
    \frac{\partial\mathcal L}{\partial s}=0
    &\iff
    -g+\tau(n-g)
    -\beta_1+\beta_2+\eta_1-\eta_2=0,
    \label{eq:stat-s}
    \\
    \frac{\partial\mathcal L}{\partial\delta}=0
    &\iff
    (1+\tau)(n-g)
    -\gamma_1+\gamma_2+\rho_1-\rho_2=0,
    \label{eq:stat-d}
    \\
    \frac{\partial\mathcal L}{\partial u_g}=0
    &\iff
    \alpha_1+\alpha_2=\lambda w_g,
    \label{eq:stat-ug}
    \\
    \frac{\partial\mathcal L}{\partial u_s}=0
    &\iff
    \beta_1+\beta_2=\lambda w_s,
    \label{eq:stat-us}
    \\
    \frac{\partial\mathcal L}{\partial u_\delta}=0
    &\iff
    \gamma_1+\gamma_2=\lambda w_\delta .
    \label{eq:stat-ud}
\end{align}


\noindent The complementarity conditions are
\begin{align}
    &\mu\,(g-n)=0,
    \label{eq:comp-mu}
    \\
    &\alpha_1\bigl(g-\widehat g-u_g\bigr)=0,
    \qquad
    \alpha_2\bigl(\widehat g-g-u_g\bigr)=0,
    \label{eq:comp-alpha}
    \\
    &\beta_1\bigl(s-\widehat s-u_s\bigr)=0,
    \qquad
    \beta_2\bigl(\widehat s-s-u_s\bigr)=0,
    \label{eq:comp-beta}
    \\
    &\gamma_1\bigl(\delta-\widehat\delta-u_\delta\bigr)=0,
    \qquad
    \gamma_2\bigl(\widehat\delta-\delta-u_\delta\bigr)=0,
    \label{eq:comp-gamma}
    \\
    &\nu_1(\underline g-g)=0,
    \qquad
    \nu_2(g-\overline g)=0,
    \label{eq:comp-nu}
    \\
    &\eta_1(\underline s-s)=0,
    \qquad
    \eta_2(s-\overline s)=0,
    \label{eq:comp-eta}
    \\
    &\rho_1(\underline\delta-\delta)=0,
    \qquad
    \rho_2(\delta-\overline\delta)=0,
    \label{eq:comp-rho}
\end{align}
together with primal feasibility and dual feasibility
\[
    \mu,\ \alpha_1,\alpha_2,\ \beta_1,\beta_2,\ \gamma_1,\gamma_2,\
    \nu_1,\nu_2,\ \eta_1,\eta_2,\ \rho_1,\rho_2\ \ge\ 0 .
\]

\noindent Based on the complementarity conditions \eqref{eq:comp-alpha} combined with primal feasibility, we know that
\begin{align}
    \label{eq:sign-info}
    \alpha_1>0 \Longrightarrow u_g=g-\widehat g \Longrightarrow g\ge\widehat g,\\
    \alpha_2>0 \Longrightarrow u_g=\widehat g-g \Longrightarrow g\le\widehat g,
\end{align}
and analogously for $(\beta_1,\beta_2)$ with $s$, and
$(\gamma_1,\gamma_2)$ with $\delta$.
\\

\noindent By \eqref{eq:stat-ug}, $\alpha_1+\alpha_2=\lambda w_g$. 
We distinguish four cases according to the values of $\alpha_1$ and $\alpha_2$.
 
\begin{itemize}
    \item \textbf{Case $\alpha_1>0$ and $\alpha_2>0$}, i.e., both epigraph constraints are active.\\
    By \eqref{eq:comp-alpha}, $u_g=g-\widehat g$ and $u_g=\widehat g-g$. 
    Subtracting these two equalities yields $2(g-\widehat g)=0$, that is, $\boxed{g=\widehat g}\in G_i(n)$.
 
    \item \textbf{Case $\alpha_1>0$ and $\alpha_2=0$.}\\
    By \eqref{eq:sign-info}, $g\ge\widehat g$, and $\alpha_1=\lambda w_g$ by \eqref{eq:stat-ug}. 
    We examine the multipliers $\mu,\nu_1,\nu_2$ of the remaining constraints on $g$ through \eqref{eq:comp-mu} and \eqref{eq:comp-nu}:
    \begin{itemize}
        \item if $\mu>0$, then $\boxed{g=n}\in G_i(n)$;
        \item if $\nu_1>0$, then $\boxed{g=\underline g}\in G_i(n)$;
        \item if $\nu_2>0$, then $\boxed{g=\overline g}\in G_i(n)$;
        \item if $\mu=\nu_1=\nu_2=0$, the stationarity condition
        \eqref{eq:stat-g} reduces to
        \begin{equation}
            \label{eq:flat-g-plus}
            \lambda w_g=-(1+\tau)(s+\delta).
        \end{equation}
        This equation \eqref{eq:flat-g-plus} constrains the prices as it forces $(s^\star,\delta^\star)$ on the line $s+\delta=-\lambda w_g/(1+\tau)$ without imposing anything on $g$.
        However, by \eqref{eq:sign-info} and primal feasibility, $g^\star\in I:=[\widehat g,\ \min\{\overline g,n\}]$ and since every $g\in I$ satisfies $g\ge\widehat g$ and $g\le n$, $t\mapsto\varphi(t,s^\star,\delta^\star)$ is affine on $I$ with slope $-(1+\tau)(s^\star+\delta^\star)-\lambda w_g$, which vanishes by \eqref{eq:flat-g-plus}. 
        The objective is therefore constant on $I$, and the optimum is attained at an endpoint of $I$, that is $\{\widehat g,\overline g,n\}\subseteq G_i(n)$.
    \end{itemize}
 
    \item \textbf{Case $\alpha_1=0$ and $\alpha_2>0$.}\\
    By \eqref{eq:sign-info}, $g\le\widehat g$, and $\alpha_2=\lambda w_g$. 
    The same three subcases $\mu>0$, $\nu_1>0$, $\nu_2>0$ give $g\in\{n,\underline g,\overline g\}\subseteq G_i(n)$. 
    If $\mu=\nu_1=\nu_2=0$, then \eqref{eq:stat-g} reduces to
    \begin{equation}
        \label{eq:flat-g-minus}
        \lambda w_g=(1+\tau)(s+\delta),
    \end{equation}
    which is again independent of $g$, this time with $g^\star\in[\underline g,\ \min\{\widehat g,n\}]$, and the same argument applies so that we can reduce the search to $\{\underline g,\widehat g,n\}\subseteq G_i(n)$.
 
    \item \textbf{Case $\alpha_1=\alpha_2=0$.}\\
    By \eqref{eq:stat-ug} this occurs only when $\lambda w_g=0$, i.e., $\lambda=0$. 
    If one of $\mu,\nu_1,\nu_2$ is positive, then $g\in\{n,\underline g,\overline g\}\subseteq G_i(n)$ as before. 
    Otherwise \eqref{eq:stat-g} gives $(1+\tau)(s+\delta)=0$.
    Since $\lambda=0$, the transportation term collapses and $t\mapsto\varphi(t,s^\star,\delta^\star)$ is affine on $[\underline g,\ \min\{\overline g,n\}]$ with slope $-(1+\tau)(s^\star+\delta^\star)=0$.
    It is therefore constant, and the optimum is also attained at an endpoint, which belongs to $\{\underline g,\overline g,n\}\subseteq G_i(n)$.
\end{itemize}

\noindent In every case we have exhibited a global maximizer of \eqref{eq:sep-minus} with $g\in G_i(n)$.
Let us move to the price variables.
This is similar to the previous case, although the interior case is slightly different.
By \eqref{eq:stat-us}, $\beta_1+\beta_2=\lambda w_s$.  

\begin{itemize}
    \item \textbf{Case $\beta_1>0$ and $\beta_2>0$.}\\
    Complementarity \eqref{eq:comp-beta} gives $u_s=s-\widehat s$ and $u_s=\widehat s-s$ so that $\boxed{s=\widehat s}\in S_i$ after subtracting.
 
    \item \textbf{Case $\beta_1>0$ and $\beta_2=0$.} Then $s\ge\widehat s$ and $\beta_1=\lambda w_s$. 
    \begin{itemize}
        \item If $\eta_1>0$ then $\boxed{s=\underline s}\in S_i$,
        \item If $\eta_2>0$ then $\boxed{s=\overline s}\in S_i$,
        \item If $\eta_1=\eta_2=0$, the stationarity condition \eqref{eq:stat-s} reduces to
        \begin{equation}
        \label{eq:flat-s-plus}
        \lambda w_s=-g^\star+\tau(n-g^\star),
    \end{equation}
    which does not depend on $s$. 
    The objective is thus affine in $s$ on the segment $[\widehat s,\overline s]$ with slope $\bigl[-g^\star+\tau(n-g^\star)\bigr]-\lambda w_s=0$ by \eqref{eq:flat-s-plus}, and we can find a global maximizer with $s\in\{\widehat s,\overline s\}\subseteq S_i$.
    \end{itemize}
 
    \item \textbf{Case $\beta_1=0$ and $\beta_2>0$.} Then $s\le\widehat s$ and $\beta_2=\lambda w_s$. 
    Similarly to the previous case, we have $s\in\{\underline s,\overline s\}$ if $\eta_1>0$ or $\eta_2>0$.
    If $\eta_1=\eta_2=0$, then \eqref{eq:stat-s} reduces to $\lambda w_s=g^\star-\tau(n-g^\star)$ and the objective is affine in $s$ on $[\underline s,\widehat s]$ with slope $\bigl[-g^\star+\tau(n-g^\star)\bigr]+\lambda w_s=0$.
    Similarly, we can find a global maximizer with $s\in\{\underline s,\widehat s\}\subseteq S_i$.
 
    \item \textbf{Case $\beta_1=\beta_2=0$}.
    This case occurs only when $\lambda w_s=0$, i.e., $\lambda=0$.
    If one of $\eta_1,\eta_2$ is positive, then $s\in\{\underline s,\overline s\}$, otherwise \eqref{eq:stat-s} gives $-g^\star+\tau(n-g^\star)=0$ and the objective is constant in $s$ on $[\underline s,\overline s]$.
    In both cases a global maximizer exists with $s\in S_i$.
\end{itemize}

\noindent The proof for $\delta$ is identical to that of $s$, using \eqref{eq:stat-ud}.
In all cases we have found a global maximizer $(g^\star,s^\star,\delta^\star)$ of \eqref{eq:deficit-epigraph} such that $g^\star\in G_i(n)$, $s^\star\in S_i$, and $\delta^\star\in\Delta_i$.
\end{proof}

\subsubsection{Under what conditions can we use strong duality?}
In our problem, we do not have (without relaxation) that the loss is convex in $\xi$, which prevents a direct use of the strong duality result for Wasserstein DRO with convex losses~\cite{EsfahaniKuhn2018}.
However, the loss is continuous and bounded on the compact support $\XiSet$, and the transportation metric is defined on a compact Polish space. 
Under these conditions, we can apply the strong duality result for Wasserstein DRO with measurable bounded losses~\cite{GaoKleywegt2022, ZhangYangGao2025}, which yields the dual representation in Proposition~\ref{prop:wasserstein-dual}.

\begin{proposition}[Wasserstein dual representation] \label{prop:wasserstein-dual} Suppose Assumptions~\ref{ass:compact-support}--\ref{ass:transportation-metric} hold. \color{red}[Carefully check this]\color{black} 
    Then, for every fixed $n\in[0,1]$, \[ \sup_{\Q\in\cP_\eps} \E_\Q[\ell(n,\xi)] = \inf_{\lambda\ge0} \left\{ \lambda\eps + \frac1N\sum_{i=1}^N \sup_{\xi\in\XiSet} \left[ \ell(n,\xi)-\lambda d_W(\xi,\hxi_i) \right] \right\}. \] 
    Consequently, the DRO problem admits the epigraphic form 
    \begin{equation} \label{eq:dro-dual-main} 
        \begin{aligned} \inf_{n,\lambda,z}\quad & \lambda\eps+\frac1N\sum_{i=1}^N z_i \\ \textnormal{s.t.}\quad & z_i\ge \mathcal V_i(n,\lambda), \qquad i=1,\ldots,N, \\ & n\in[0,1], \qquad \lambda\ge0, 
        \end{aligned} 
    \end{equation} where 
    \begin{equation} \label{eq:separation-problem-dual} \mathcal V_i(n,\lambda) := \sup_{\xi\in\XiSet} \left\{ \ell(n,\xi)-\lambda d_W(\xi,\hxi_i) \right\}. \end{equation} 
\end{proposition} 
\begin{proof} For every fixed $n$, the loss is continuous on the compact set $\XiSet$, hence bounded and Borel measurable. 
    The metric space $(\XiSet,d_W)$ is compact, and therefore Polish. 
    The result follows from the strong duality theorem for Wasserstein distributionally robust optimization with measurable bounded losses~\cite{GaoKleywegt2022, ZhangYangGao2025}. 
    Since the reference distribution is empirical, the dual expectation reduces to a finite average over the observed samples. 
\end{proof} 

\color{red}See \cite{GaoKleywegt2022, ZhangYangGao2025} for the general result, precise assumptions and proof.
\color{black}


\begin{remark}
The exact finite separation can also be used inside a cutting-plane or outer-approximation algorithm: given a candidate pair $(n,\lambda)$, compute $\mathcal V_i(n,\lambda)$ exactly by \eqref{eq:finite-separation}, add the corresponding violated epigraph inequality, and iterate. 
\end{remark}

\subsubsection{Critical Wasserstein penalties} 
\label{subsec:critical-lambda} 
The dual multiplier $\lambda$ has a useful interpretation. 
It is the marginal cost assigned by the decision maker to transporting probability mass away from the empirical samples. Small values of $\lambda$ allow the adversary to move samples substantially inside the support, leading to robust behavior close to worst-case optimization. 
Large values of $\lambda$ strongly penalize transportation, forcing the adversary to remain close to the empirical distribution and recovering behavior close to SAA. 
For a fixed $n$, define the empirical and robust values \[ F_{\rm SAA}(n) = \frac1N\sum_{i=1}^N \ell(n,\hxi_i), \qquad F_{\rm RO}(n) = \sup_{\xi\in\XiSet}\ell(n,\xi). \] 
The Wasserstein DRO value satisfies \[ F_{\eps}(n) = \inf_{\lambda\ge0} \left\{ \lambda\eps + \frac1N\sum_{i=1}^N \mathcal V_i(n,\lambda) \right\}. \] 
When $\eps=0$, the optimal value reduces to $F_{\rm SAA}(n)$. 
When $\eps$ is large enough, the transportation budget ceases to be the binding limitation, and the adversarial distribution may concentrate on worst-case support points, producing the robust value $F_{\rm RO}(n)$. 
The transition \[ \mathrm{SAA} \longrightarrow \mathrm{DRO} \longrightarrow \mathrm{RO} \] depends jointly on the radius $\eps$, the support $\XiSet$, the metric weights $W$, and the magnitude of the adverse price-generation interactions in $\ell$. 

\paragraph{SAA $\to$ DRO transition}
Consider the Lipschitz constant of the loss with respect to the transportation metric:
\begin{equation} 
    \label{eq:lipschitz-threshold} L(n) := \sup_{\xi\neq\xi'\in\XiSet} \frac{|\ell(n,\xi)-\ell(n,\xi')|} {d_W(\xi,\xi')}. 
\end{equation} 
We know that for any $\lambda\ge L(n)$, \color{red}[Insert proof here]\color{black} the transportation penalty dominates the separation problem and prevents adversarial movement, so that \[ \sup_{\xi\in\XiSet} \left\{ \ell(n,\xi)-\lambda d_W(\xi,\hxi_i) \right\} = \ell(n,\hxi_i). \] 
Thus, for a sufficiently large $\lambda$, we recover the empirical loss, and the critical penalty $\lambda=L(n)$ marks the transition from SAA to DRO.


\paragraph{DRO $\to$ RO transition}
Conversely, when $\lambda$ is small, the transportation penalty is weak and the separation problem moves toward support points that maximize the loss. 
We recover the robust optimization value.
\color{red}[How to find this critical $\lambda$? Bissection algorithm?]\color{black}

\section{Results}
\subsection{Experimental design}
\subsubsection{Choice of params (cf empirical analysis)}
\subsubsection{Wasserstein radius automatic calibration}
\subsubsection{Weight calibration}
\subsubsection{Out-of-sample evaluation}
\subsection{Performance metrics}
\subsubsection{Mean out-of-sample profit}
\subsubsection{CVaR}
\subsubsection{Reliability}
\subsubsection{Conservatism gap}
\subsubsection{Runtime}
\subsubsection{Nomination stability}
\subsection{Positive-prices scenario (synthetic data)}
\subsubsection{Baseline scenario}
\subsubsection{Tight support}
\subsubsection{Wide support}
\subsubsection{Heavy-tailed imbalance prices}
\subsubsection{Rescaled uncertainties}
\subsubsection{Non-box support}
\subsection{Negative-prices scenario (synthetic data)}
\subsubsection{Baseline scenario}
\subsubsection{Tight support}
\subsubsection{Wide support}
\subsubsection{Heavy-tailed imbalance prices}
\subsubsection{Rescaled uncertainties}
\subsubsection{Non-box support}
\subsection{Positive-prices scenario (actual data)}
\subsubsection{Baseline scenario}
\subsubsection{Tight support}
\subsubsection{Wide support}
\subsubsection{Heavy-tailed imbalance prices}
\subsubsection{Rescaled uncertainties}
\subsubsection{Non-box support}
\subsection{Negative-prices scenario (actual data)}
\subsubsection{Baseline scenario}
\subsubsection{Tight support}
\subsubsection{Wide support}
\subsubsection{Heavy-tailed imbalance prices}
\subsubsection{Rescaled uncertainties}
\subsubsection{Non-box support}
\subsection{Discussion}
\subsubsection{Synthetic vs actual data}
\subsubsection{Exact vs approx}
\subsubsection{Box vs non-box}
\subsubsection{Impact of negative prices}
\subsubsection{Impact of extreme events}
\subsubsection{Impact of scaling}
\subsubsection{Limitations}
\subsubsection{Future work}


\begin{thebibliography}{9}

\bibitem{EsfahaniKuhn2018}
P. M. Esfahani and D. Kuhn. 2018.
Data-driven distributionally robust optimization using the Wasserstein metric: performance guarantees and tractable reformulations.
\emph{Mathematical Programming}.

\bibitem{HomemDeMello}
C. A. Gamboa, D. M. Valladao, A. Street, and T. Homem-de-Mello. 2021. Decomposition methods for Wasserstein-based data-driven distributionally robust problems
\emph{Operations Research Letters}.

\bibitem{AolariteiBangouraBolognaniLanzettiDorfler2025} 
Aolaritei, L., Bangoura, B., Bolognani, S., Lanzetti, N., and Dörfler, F. (2025). 
\textit{Hedging against Black Swans in Day-Ahead Energy Markets}.

\bibitem{GaoKleywegt2022} 
Gao, Rui and Kleywegt, Anton J. (2022). 
\textit{Distributionally Robust Stochastic Optimization with Wasserstein Distance}.

\bibitem{ZhangYangGao2025} 
Zhang, Luhao and Yang, Jincheng and Gao, Rui (2025). 
\textit{A Short and General Duality Proof for Wasserstein Distributionally Robust Optimization}.
\end{thebibliography}


\end{document}
